\section{Theory}
\label{sec:theory}

\subsection{The price of a budget, and two floors}
\label{sec:theory-main}
Theorem~\ref{thm:main} takes conditionally i.i.d.\ views with lag tail $G$, and $d$ is the binary relative entropy. In (b) and (c) the policy is
\emph{verification-only}, ending by retrying or declaring, with $\PP(\text{declare}\mid X{=}0)\le\beta$. In (b) it decides by $r_e{+}w$ from end-to-end reads alone,
and in (c) it meets $(\alpha,\beta)$ with $\alpha{+}\beta{<}1$ and finite $\E_x[N]$. In (d) it never asks, the effect commits with prior $\pi$, and $\bar\pi=\min(\pi,1{-}\pi)$ is the prior floor.

\begin{theorem}
\label{thm:main}
(a)~If $G(w){<}\alpha$, waiting until $r_e{+}w$ and then an SPRT meet $(\alpha,\beta)$. (b)~$\PP(\text{retry}\mid X{=}1)\ge(1{-}\beta)G(w)$.
(c)~$\E_0[N]\ge d(1{-}\beta\,\|\,\alpha)/D_0$. (d)~$\PP(\dup){+}\PP(\loss)\ge\bar\pi\big(1{-}\sqrt{D_0\E_0[N]/2}\big)$.
\end{theorem}

None of these inequalities is new, and Appendix~\ref{app:proofs} states them in full. Part~(a) is Wald's bound~\cite{wald1945sprt} plus a union bound on
the lag, and (b) is the noisy-read, finite-wait form of Li's waiting argument~\cite{li2026exactlyonce}. Part~(c) says a verification-only policy needs about $\log(1/\alpha)/D_0$ reads
(Lemma~1 of~\cite{kaufmann2016complexity}, see also~\cite{tartakovsky2014sequential}). Part~(d) says a policy that never asks errs at least at the Bayes floor, by Le Cam's
two-point method~\cite{tsybakov2009nonparametric}.

\begin{corollary}
\label{cor:notfree}
If $0<\epsilon_1<1$, or if $G(t)>0$ for every $t$, every verification-only policy with $\PP(\text{retry}\mid X{=}0)>0$ has $\PP(\text{retry}\mid X{=}1)>0$: $\alpha=0$ forces $\beta=1$,
for unbounded stopping times too. For $\alpha,\beta>0$ with $\alpha+\beta<1$, $(\alpha,\beta)$ is reachable by verification if and only if $D_0>0$.
\end{corollary}

Exactly-once by verification is a limit, not a point, so \sys exposes $(\alpha,\beta)$ as a knob. By~(b), asking or a lag-free family such as the app's
Outbox is the only escape from the lag floor. Otherwise the lag share is paid in waiting time (Table~\ref{tab:lag}), and the planner asks when the certified wait exceeds a cap. By~(d), on a look-alike screen ($D_0\to0$) the error approaches the prior floor for any number of reads.

\subsection{Correlated reads: certify views, not reads}
\label{sec:theory-corr}
Theorem~\ref{thm:main} counts every read as fresh evidence. A VLM re-reading a screenshot does not supply that, and neither does an OCR pass over a hidden nonce field.
\alg therefore reads each view once and certifies \emph{family verdicts}. A family $f$ returns the majority side $s\in\{\mathrm{abs},\mathrm{com}\}$ and the number $r$ of its views
on that side, or nothing on a tie. Let $q^{\mathrm{abs}}_f(r)$ bound the probability that at least $r$ views of $f$ read absent when $X{=}1$ and the effect is visible. Let
$q^{\mathrm{com}}_f(r)$ bound the probability that at least $r$ read committed when $X{=}0$.

Product certification (Proposition~\ref{prop:family} in the appendix) consults families in a fixed order and multiplies the bounds of the verdicts seen so far. Under A2 the
probability of reaching a retry is at most a sum $B(\tau)$ of such products, with no assumption about the views within a family. The planner picks the largest $\tau$ with
$B(\tau)\le\alpha_{\mathrm{test}}$. If a hop-1 entry ends failed or stuck, the planner hands over once to the end-to-end families, and the same bound covers that path.
The lag adds two terms. Let $w_1$ be the wait after a hop-1 ``absent'' verdict and $w_0$ the wait on every other path (hop~1 tied or absent from the app, an entry of an app without idempotent retries,
a failed or stuck entry). Both waits run from the anchor $r_e$. For an effect of class $c$,
\begin{equation}
\PP(\text{retry}\mid X{=}1)\le B(\tau)+q_h\,G_c(w_1)+G_c(w_0),
\label{eq:lagsplit}
\end{equation}
where $q_h$ bounds the hop-1 ``absent'' verdict. The first lag term uses only that hop-1 readings are independent of $T$ given $X$, and the second needs no independence
(Appendix~\ref{app:proofs}). \alg certifies each term at $\alpha_{\mathrm{lag}}/2$, or all of
$\alpha_{\mathrm{lag}}$ to $w_0$ if the app has no hop-1 family.

A2 lets bounds multiply, and it is doubtful when families read the same screens: a stale history after a restart fools nonce, digest and VLM at
once~\cite{zheng2026vpcontrol}. The remedy is to calibrate larger units.

\begin{proposition}[Joint certification]
\label{prop:joint}
Fix a score that orders verdict tuples (products of bounds from a pilot split) and let $R_\tau$ be the set of tuples on which the cascade retries at threshold $\tau$, united
over the configurations the planner runs (with or without hop~1 and the VLM). Test $\tau_1<\tau_2<\cdots$ in order on $n$ fresh calibration effects and keep the largest
$\tau_j$ such that, for all $i\le j$, the one-sided $(1-\delta)$ CP bound on $\PP(R_{\tau_i}\mid X{=}1)$ is at most $\alpha_{\mathrm{test}}$. With probability at least $1-\delta$,
$\PP(\text{retry}\mid X{=}1,\text{visible})\le\alpha_{\mathrm{test}}$, whatever the dependence across families. Certifying needs about $n\ge\ln(1/\delta)/\alpha_{\mathrm{test}}$ effects.
\end{proposition}
The pilot that fixes the score is disjoint from the $n$ test effects. \sys uses the first tenth of each outcome's calibration effects as the pilot and tests on the other nine tenths,
so $n$ counts only the latter. An intermediate mode, source families, keeps A2 only between hop~1 and the history readers, two \emph{evidence sources} (Appendix~\ref{app:proofs}; Element fails
its check). One family cannot certify a side below its persistent blind mass however many views it has (Lemma~\ref{lem:floor} in the appendix).
A hidden nonce field is such a blind state, so certifying ``absent'' below its rate needs a second family.

\subsection{Calibration}
\label{sec:theory-cal}
\begin{lemma}[Conservative calibration]
\label{lem:calib}
(i)~For a family of i.i.d.\ binary views, running the SPRT with $\tilde\epsilon_0\ge\epsilon_0$, $\tilde\epsilon_1\ge\epsilon_1$ ($\tilde\epsilon_0+\tilde\epsilon_1<1$) preserves the
$(\alpha,\beta)$ guarantee.
(ii)~One fault class gives an effect's decision $K$ bounds: verdict levels or regions ($q_h$ is the largest hop-1 ``absent'' level) and the two waits. Suppose each is a one-sided $(1-\delta)$ Clopper-Pearson
bound~\cite{clopper1934use} from calibration effects that are i.i.d.\ draws from the population of interrupted deployment effects of the effect's observable fault
class. Suppose also that an effect of unknown class uses the maximum over classes. Then with probability
at least $1-Kn_{\mathrm{cls}}\delta$ all hold, where $n_{\mathrm{cls}}$ is the number of fault classes. Marginally, $\PP(\text{retry}\mid X{=}1)\le\alpha_{\mathrm{test}}+\alpha_{\mathrm{lag}}+Kn_{\mathrm{cls}}\delta$ and $\PP(\text{declare}\mid X{=}0)\le\beta_{\mathrm{test}}+Kn_{\mathrm{cls}}\delta$.
\end{lemma}
Part~(i) is a monotone coupling and (ii) a union bound. Keying by class makes the lab's fault mix irrelevant at the price of labels per class. The unknown-class fallback
uses all classes' bounds, so the union runs over $Kn_{\mathrm{cls}}$ bounds. \sys sets $\delta=\min(\alpha,\beta)/(10\,K\,n_{\mathrm{cls}})$, with $n_{\mathrm{cls}}$ the number of fault classes when bounds are keyed by class (4 in the simulations, 3 in the held-out runs) and $n_{\mathrm{cls}}{=}1$ when they are pooled.
Here $K{=}11$ for product certification (five ``absent'' and four ``committed'' verdict levels, two waits; the nonce's false-match bound is analytic), $K{=}8$ for source
families and $K{=}5$ for joint certification. Then $Kn_{\mathrm{cls}}\delta=\min(\alpha,\beta)/10\le\min(\alpha_{\mathrm{cal}},\beta_{\mathrm{cal}})$, and the guarantee pays the union bound on both sides.
The population assumption is the one that bites: labels from uninterrupted effects miss cached history after a kill, optimistic ``sending'' entries and blackout lags.
Lemma~\ref{lem:drift} in the appendix bounds drift that is independent across families. Common-mode drift is what joint certification is for (\S\ref{sec:eval-pop}).

\subsection{Mobile structure: nonces and two hops}
\label{sec:twohop}
A content-carried nonce turns look-alike effects into text the probe can match. The planner writes $k$ random symbols from an alphabet of size $m$ and reuses them on retry.
If the probe never reads drafts, failed entries or notifications, a view with $M$ visible tokens and match distance $t$ has false-match rate at most $M|B_t|/m^k$, whatever
look-alikes are on screen, where $|B_t|$ counts the strings near a token. Its miss rate is at most the hidden-field rate plus OCR errors (Lemma~\ref{lem:nonce}).
Reusing the nonce also makes a later duplicate detectable, because two items carry the same tag.

\begin{theorem}[Two-hop decomposition]
\label{thm:twohop}
Let hop~1 create $n_1$ entries in the app's local outbox and hop~2 transmit entries to the server. Assume three conditions.
(H1)~The server is reached only through the outbox.
(H2)~Each entry has an app-assigned id, unique per entry, persisted with it and reused by every retransmission from the same device, also after restarts.
(H3)~Admissible fault schedules have finitely many agent kills, app kills that do not cancel the app's scheduled work, transmissions lost before the server applies them,
and acknowledgements lost after. Every entry is transmitted at least once on the fault-free schedule, and a positive window separates an entry's creation from its first
transmission. Every retransmission happens within the server's \emph{dedup horizon} $H_{\mathrm{dedup}}$ of the first application.
Let \textbf{G2} mean: entries are durable, each is retransmitted until acknowledged, and the server applies each (device, id) at most once within $H_{\mathrm{dedup}}$. Then $C_e=1$ under every admissible
schedule if and only if $n_1=1$ and G2 holds. If retransmission may instead stop in a failed state shown by the app (\textbf{G2$'$}), then $n_1=1$ implies $C_e\le1$ under
every admissible schedule, and $C_e=0$ only if the entry ends failed.
\end{theorem}
The positive window in H3 is a property of the app. Element satisfies H3, but its window is a fraction of a second, shorter than one accessibility read. So no probe sees
an Element entry before it is sent, and Element has no hop-1 family (\S\ref{sec:eval-real}).
($\Leftarrow$) is at-least-once delivery plus idempotent application~\cite{helland2012idempotence,birrell1984rpc,wang2021kafka}. What the agent gains is that hop~1 is
observable on the device within seconds, so for an app with idempotent retries the late-commit regime of Theorem~\ref{thm:main}(b) becomes a local check. The planner therefore counts hop-1
entries created after $t_0$ and hands a pending entry over to the end-to-end path before its age approaches $H_{\mathrm{dedup}}$ (\S\ref{sec:design-planner}). ``Retrying is safe'' thus means the app's own retry, on the same device, within $H_{\mathrm{dedup}}$.
An agent that redoes the step through the GUI creates a new entry with a new id ($n_1{=}2$), which no server merges, and a re-login creates a new device. A failed entry is
ambiguous end to end, because the server may have applied it before the app gave up, so it goes to the end-to-end path. K-9~Mail enters this state by design after every lost
SMTP reply (\S\ref{sec:measure-teaser}). Force-stop cancels scheduled work and violates H3, so a force-stopped app is treated as
an app without idempotent retries until its entries reach a terminal state. This only removes the hop-1 shortcut. Element and DAVx5 still resolved force-stops, outside the theorem (Appendix~\ref{app:proofs}, with dedup scopes per protocol). Per-step guarantees compose across a workflow by a union bound
(Proposition~\ref{prop:saga}).
